% Fichamento — Desenvolvimento de Software Suportado por IA
% Mestrado em Computação Aplicada · IFES Serra · 2026.2
% Uma página por artigo. Compile com pdflatex.

\documentclass[11pt,a4paper]{article}

\usepackage[utf8]{inputenc}
\usepackage[T1]{fontenc}
\usepackage[brazil]{babel}
\usepackage[margin=2.5cm]{geometry}
\usepackage{hyperref}
\usepackage{titlesec}
\usepackage{parskip}
\usepackage{csquotes}

\hypersetup{colorlinks=true, urlcolor=black, linkcolor=black}
\titleformat{\section}{\bfseries\large}{}{0pt}{}
\titlespacing*{\section}{0pt}{12pt}{2pt}
\pagestyle{empty}

% ---------------------------------------------------------------
% Preencha os campos abaixo
% ---------------------------------------------------------------
\newcommand{\aluno}{Nome do aluno}
\newcommand{\datafichamento}{\today}

\begin{document}

{\Large\bfseries Fichamento}\\[2pt]
{\small \aluno{} \quad·\quad \datafichamento}

\hrule\vspace{6pt}

\section*{Referência}
Autor, título, ano e link.\\
Ex.: FAN, A. et al. \emph{Large Language Models for Software Engineering:
Survey and Open Problems}. ICSE-FoSE, 2023.
Disponível em: \url{https://arxiv.org/abs/2310.03533}.

\section*{Tema}
Assunto principal do texto.

\section*{Objetivo}
O que o autor pretende discutir ou demonstrar.

\section*{Resumo}
Síntese das principais ideias, com suas próprias palavras.

\section*{Citação importante}
\begin{quote}
\enquote{Trecho relevante do texto.} (p.~XX)
\end{quote}

\section*{Comentário}
Sua avaliação sobre o texto e como ele pode ser útil para sua pesquisa.

\end{document}
